% !PRINTED 24
\weberpage{24}
We again introduce homogeneous coordinates $x_1$, $x_2$, $x_3$, $x_4$ and
take the equations of two given surfaces of the second degree in
the form:
\weq{3}{\begin{aligned} \varphi(x_1,x_2,x_3,x_4) &= \Sigma a_{ik}x_ix_k, \\ \psi(x_1,x_2,x_3,x_4) &= \Sigma b_{ik}x_ix_k. \end{aligned}}

The two surfaces determine a \gesperrt{pencil of surfaces} of the second
order $\xi\varphi=\eta\psi$. All surfaces of the pencil intersect in a space
curve of the fourth order, which we call the \gesperrt{fundamental curve}.
One obtains the same curve and the same pencil if the functions
$\varphi$ and $\psi$ are replaced by $\varphi'$, $\psi'$, which are derived from $\varphi$, $\psi$
by means of the linear substitution $\begin{pmatrix} m & n \\ p & q \end{pmatrix}$, that is, by
\weq{4}{\begin{aligned} \varphi' &= m\varphi+n\psi, \\ \psi' &= p\varphi+q\psi \end{aligned}}
and whose determinant $r=mq-np$ is different from zero. From
this there results the identity
\weq{5}{\xi\varphi+\eta\psi = \xi'\varphi'+\eta'\psi',}
in which
\weq{6}{\begin{aligned} \xi &= m\xi'+p\eta', \\ \eta &= n\xi'+q\eta' \end{aligned}}
represents a linear transformation of the variables $\xi$, $\eta$. Besides
this \gesperrt{binary substitution} there also comes into consideration a linear quaternary substitution of the variables $x_1$, $x_2$, $x_3$, $x_4$ (coordinate transformation), whose determinant we may assume $=1$.

Every quadratic form possesses, with respect to these latter transformations, an \gesperrt{invariant}, namely the \gesperrt{Hesse
determinant} (Bd.\,I, § 62, 63, 66), and we thus obtain as the
invariant of the form (5):
\weq{7}{f(\xi,\eta) = \Sigma \pm(\xi a_{11}+\eta b_{11})(\xi a_{22}+\eta b_{22})(\xi a_{33}+\eta b_{33})(\xi a_{44}+\eta b_{44}).}

This is a binary biquadratic form in the variables $\xi$, $\eta$,
which, developed, we represent thus:
\weq{8}{f(\xi,\eta) = a_0\xi^4+a_1\xi^3\eta+a_2\xi^2\eta^2+a_3\xi\eta^3+a_4\eta^4.}

The coefficients of this biquadratic form $a_0$, $a_1$, $a_2$, $a_3$, $a_4$
do not change under a coordinate transformation. They are
therefore called \gesperrt{simultaneous invariants} of the pair of forms $\varphi$, $\psi$ ($a_0$ and
$a_4$ are the determinants of $\varphi$ and of $\psi$).
